% =====================================================================
%  Thème 5 — EXERCICES — Niveau DIFFICILE
% =====================================================================
\documentclass[11pt,a4paper]{article}
\usepackage[utf8]{inputenc}
\usepackage[T1]{fontenc}
\usepackage{lmodern}
\usepackage[french]{babel}
\usepackage{amsmath,amssymb,mathtools}
\usepackage{icomma}
\usepackage{xcolor}
\usepackage[most]{tcolorbox}
\usepackage{enumitem}
\usepackage{array}
\usepackage{booktabs}
\usepackage{fancyhdr}
\usepackage[hidelinks]{hyperref}
\usepackage[a4paper,margin=2.2cm,headheight=26pt,headsep=10pt]{geometry}

\definecolor{primary}{RGB}{223,87,99}
\definecolor{primarydark}{RGB}{150,42,55}

\tcbset{boxbase/.style={enhanced,breakable,boxrule=0.9pt,arc=2.5pt,
   left=3mm,right=3mm,top=2.3mm,bottom=2.3mm,
   fonttitle=\bfseries,coltitle=white,toptitle=0.6mm,bottomtitle=0.6mm}}
\newtcolorbox[auto counter]{exercice}[1][]{boxbase,colback=primary!4,colframe=primary,
  title={\textsc{Exercice}~\thetcbcounter\ifx\relax#1\relax\relax\else\textmd{ --- \itshape #1}\fi}}

\pagestyle{fancy}
\fancyhf{}
\fancyhead[L]{\small\textcolor{primarydark}{\textbf{Fiche 6} — Les limites}}
\fancyhead[R]{\small\textcolor{primarydark}{\textsc{Thème 5 — Difficile}}}
\fancyfoot[C]{\small\thepage}
\renewcommand{\headrulewidth}{0.8pt}
\renewcommand{\headrule}{\hbox to\headwidth{\color{primary}\leaders\hrule height \headrulewidth\hfill}}

\newcommand{\R}{\mathbb{R}}

\begin{document}

\begin{center}
{\Large\bfseries\color{primarydark}Thème 5 — L'Hospital \& limites trigonométriques}\\[2pt]
{\large\color{primary}Exercices — Niveau Difficile}
\end{center}
\vspace{4pt}
{\small\itshape Niveau concours : combinaison fine de limites remarquables, deux méthodes pour un même
résultat, et une mise en situation avec l'Hospital répété.}
\vspace{6pt}

\begin{exercice}[combiner trois limites de référence]
On veut $\displaystyle\lim_{x\to 0}\frac{\tan x-\sin x}{x^3}$.
\begin{enumerate}[label=\alph*),leftmargin=2em,itemsep=3pt]
  \item Montrer que $\tan x-\sin x=\dfrac{\sin x\,(1-\cos x)}{\cos x}$.
  \item En déduire l'écriture $\dfrac{\tan x-\sin x}{x^3}=\dfrac{\sin x}{x}\cdot\dfrac{1-\cos x}{x^2}\cdot\dfrac{1}{\cos x}$.
  \item Sachant $\dfrac{\sin x}{x}\to 1$, $\dfrac{1-\cos x}{x^2}\to\dfrac12$ et $\dfrac{1}{\cos x}\to 1$,
        conclure la valeur de la limite.
\end{enumerate}
\end{exercice}

\begin{exercice}[deux méthodes pour un quotient de sinus]
On veut $\displaystyle\lim_{x\to 0}\frac{\sin 5x}{\sin 3x}$.
\begin{enumerate}[label=\alph*),leftmargin=2em,itemsep=3pt]
  \item Vérifier que la forme est $\frac00$.
  \item \textbf{Méthode 1 (limites remarquables) :} écrire
        $\dfrac{\sin 5x}{\sin 3x}=\dfrac{\sin 5x}{5x}\cdot\dfrac{3x}{\sin 3x}\cdot\dfrac{5x}{3x}$ et conclure.
  \item \textbf{Méthode 2 (l'Hospital) :} dériver haut et bas, puis conclure. Vérifier qu'on retrouve le
        même résultat.
  \item Généraliser : que vaut $\displaystyle\lim_{x\to 0}\frac{\sin(ax)}{\sin(bx)}$ pour $a,b\neq 0$ ?
\end{enumerate}
\end{exercice}

\begin{exercice}[mise en situation : erreur de l'approximation des petits angles]
En physique, pour de petits angles $\theta$ (en radians), on approxime $\sin\theta\approx\theta$
(pendule, optique\dots). On étudie l'erreur relative via
\[ E(\theta)=\frac{\theta-\sin\theta}{\theta^3}. \]
\begin{enumerate}[label=\alph*),leftmargin=2em,itemsep=3pt]
  \item Justifier d'abord que $\displaystyle\lim_{\theta\to 0}\frac{\sin\theta}{\theta}=1$ (ce qui légitime $\sin\theta\approx\theta$).
  \item Montrer que $E(\theta)$ est de la forme $\frac00$, puis appliquer la règle de l'Hospital
        \emph{trois fois} pour calculer $\displaystyle\lim_{\theta\to 0} E(\theta)$.
  \item En déduire que, pour $\theta$ petit, $\theta-\sin\theta\approx \dfrac{\theta^3}{6}$. L'erreur
        commise en écrivant $\sin\theta\approx\theta$ est-elle « petite » devant $\theta$ ? Justifier.
\end{enumerate}
\end{exercice}

\end{document}
